% Chapter 1 -- Introduction (O-RAN / BMPP-DQN track)
%
% Governs: manuscript/ORAN_BMPP_DQN_Concept_Note_v1.md (Sections 2-4, 6, 7)
% Chapter mapping: docs/oran_thesis_guide.md
% Status: DRAFT. Intended to be \input{} from the thesis-wide main.tex once
% that skeleton exists; a thesis-wide references.bib does not yet exist, so
% citations below are numbered locally (IEEE style) against the short
% reference list at the end of this file. Renumber/merge into the
% thesis-wide bibliography once main.tex is created.
%
% Distinguishes this file from any future C-RAN-track "Chapter 1" via the
% "oran" filename marker, per the repository's own oran_* naming convention
% (AGENTS.md).

\chapter{Introduction}
\label{chap:oran-introduction}

\section{Background}
\label{sec:oran-background}

The transition from monolithic base station architectures toward Open Radio
Access Networks (O-RAN) is one of the defining architectural shifts of 5G
and its evolution toward 6G. O-RAN disaggregates the traditional gNB into
three functional entities -- the Radio Unit (RU), Distributed Unit (DU), and
Central Unit (CU) -- connected by open, standardized fronthaul and midhaul
interfaces rather than proprietary vendor links~\cite{oran2021mvp}. This
disaggregation, together with the introduction of RAN Intelligent
Controllers (RICs), is intended to make the RAN programmable: a Non-Real-Time
RIC operates on a $1$--$10$\,s control loop for slower policy and
configuration decisions, while a Near-Real-Time RIC operates on a
$10$--$100$\,ms loop for faster radio-resource decisions~\cite{oran2021mvp}.

This programmability comes at the cost of substantially increased
operational complexity. The RAN segment alone is estimated to account for
$70$--$80\%$ of a mobile network's total energy consumption, making RAN
energy efficiency a first-order operating-expenditure concern for operators
and a first-order sustainability concern for the industry. Unlike a
monolithic base station, an O-RAN deployment exposes energy-relevant control
knobs at every disaggregation boundary: whether an RU is powered on or
placed into a low-power state, which functional split is used to divide
processing between the RU and the DU/CU, how much transmit power and how
large a share of physical resource blocks (PRBs) each RU is allocated. These
decisions are naturally structured as a \emph{hybrid} action space: some are
discrete (RU on/off, split selection) and some are continuous (transmit
power, PRB share), and they occur at different natural timescales -- RU
activation and split selection are Non-RT-RIC-appropriate, slow decisions,
while power and PRB allocation are Near-RT-RIC-appropriate, fast decisions.

Deep Reinforcement Learning (DRL) is a natural candidate for this class of
sequential decision problem, since it can learn a control policy directly
from interaction with a (simulated or real) environment without requiring an
explicit, tractable model of the channel and traffic dynamics. However, the
specific structure of the O-RAN energy-optimization problem -- a hybrid,
multi-timescale, discrete-continuous action space -- does not fit any single
standard DRL algorithm cleanly. Deep Q-Networks (DQN) and their variants
handle discrete actions only; Deep Deterministic Policy Gradient (DDPG) and
its successors handle continuous actions only. Parameterized-action methods
such as P-DQN~\cite{xiong2018pdqn} and its refinement, Multi-Pass DQN
(MP-DQN)~\cite{bester2019mpdqn}, were designed specifically to couple a
discrete action choice with a continuous parameter for that action, but both
process the full joint action space in a single, flat network, so one
decision's learning signal can interfere with another's. Branching DQN
(BDQ)~\cite{tavakoli2018branching} instead decomposes a discrete action
space into $R$ independent branches, reducing the size of the network's
output from $\prod_r |\mathcal{A}_r|$ to $\sum_r |\mathcal{A}_r|$, but is
limited to purely discrete action spaces. No existing algorithm combines
branching's discrete-action decomposition with a parameterized method's
principled handling of coupled discrete-continuous actions.

\section{Problem Statement}
\label{sec:oran-problem-statement}

The most recent published O-RAN energy-optimization DRL framework, OREO
(Open RAN Energy Optimization)~\cite{qazzaz2026oreo}, illustrates the
consequence of this gap. OREO uses Proximal Policy Optimization (PPO) with a
purely continuous action representation, recovering discrete RU on/off
decisions by thresholding a continuous output -- an approximation that
destroys the gradient signal a genuinely discrete decision would otherwise
provide, rather than a principled treatment of the hybrid action space. Its
action representation is flat, giving no mechanism to prevent one decision
type from interfering with another's learning, its reward weighting is
static rather than adaptive to changing network conditions, and its control
loop operates at a single timescale (Non-RT RIC only), leaving the faster
Near-RT RIC decisions -- and any coupling between the two timescales --
outside its scope.

This motivates the central problem addressed by this thesis: \emph{how can a
DRL framework combine action-space branching and multi-pass parameterized
action processing to efficiently optimize energy consumption in O-RAN, while
correctly handling the hybrid discrete-continuous action space that arises
from RU activation, functional split selection, transmit power control, and
PRB allocation under dynamic traffic conditions?} This thesis proposes and
evaluates \emph{BMPP-DQN} (Branching Multi-Pass Parameterized DQN), a DRL
architecture designed specifically to close this gap.

\section{Research Questions}
\label{sec:oran-research-questions}

This thesis is organized around three research questions:

\begin{description}
  \item[RQ1] How can branching action decomposition and multi-pass
    parameterized processing be integrated into a single DRL architecture
    to efficiently handle O-RAN's hybrid discrete-continuous action space?
  \item[RQ2] What is the energy-saving performance of BMPP-DQN relative to
    standard DQN, DDPG, and MP-DQN baselines in a simulated O-RAN
    environment?
  \item[RQ3] How does explicitly separating slow (RU activation, split
    selection) and fast (power, PRB allocation) decisions across two
    timescales affect learning convergence and the resulting energy
    efficiency, relative to treating all decisions at a single timescale?
\end{description}

\section{Research Objectives}
\label{sec:oran-objectives}

The corresponding research objectives are to:

\begin{enumerate}
  \item Design and implement the BMPP-DQN algorithm, integrating a
    branching discrete-action architecture with multi-pass parameterized
    continuous-action processing within a single coupled network.
  \item Formulate O-RAN energy optimization as a parameterized Markov
    Decision Process (MDP) with four action branches -- RU activation,
    functional split selection, transmit power, and PRB allocation fraction
    -- spanning two decision timescales.
  \item Implement a focused, single-gNB O-RAN simulation environment with
    a tractable, three-option abstraction of 3GPP functional splits.
  \item Quantify BMPP-DQN's energy-saving performance against DQN, DDPG,
    and MP-DQN baselines under matched traffic and evaluation conditions,
    against a target of $\geq 15\%$ energy savings relative to these
    baselines.
  \item Analyze how the two-timescale, branch-separated design affects
    learning convergence and overall energy efficiency, relative to a
    single-timescale treatment of the same action space.
\end{enumerate}

\section{Significance of the Research}
\label{sec:oran-significance}

This work contributes to the DRL and O-RAN energy-optimization literatures
in three ways. First, it integrates branching decomposition and multi-pass
parameterized processing into a single coupled architecture -- a combination
that, to the best of the author's knowledge, has not previously been
proposed -- extending the DRL literature on hybrid discrete-continuous
action spaces. Second, it targets a specific, practically important
application: O-RAN's disaggregated RU/DU/CU architecture directly exposes
the discrete-continuous, multi-timescale control problem this design is
built for, and RAN energy consumption is a direct, ongoing operating-expense
and sustainability concern for network operators. Third, by comparing
BMPP-DQN against three representative baselines (a purely discrete method,
a purely continuous method, and an existing parameterized method) under a
common simulation environment, this thesis provides an incremental,
evidence-based advancement over existing parameterized DRL methods, rather
than an untested architectural proposal.

\section{Scope and Limitations}
\label{sec:oran-scope}

The scope of this thesis is deliberately bounded to keep the research
question tractable within an MPhil timeline. The study considers a single
gNB, modeled with multiple RUs, one DU, and one CU, and a tractable
three-level abstraction of 3GPP functional split options (approximating
Options 2, 6, and 8 of 3GPP TR~38.801) rather than the full space of
standardized splits. Only downlink energy optimization is considered.
BMPP-DQN is compared against three baselines -- DQN, DDPG, and MP-DQN --
rather than the full space of prior hybrid-action methods; OREO is discussed
qualitatively (Section~\ref{sec:oran-problem-statement}) but not reproduced,
given the complexity of its PPO-based training pipeline. All evaluation is
simulation-based, using a simplified wireless channel model and a
literature-informed but not independently validated power-consumption model
(discussed further, with its supporting evidence, in
Chapter~\ref{chap:oran-system-model}). No real-world O-RAN testbed
validation, multi-gNB or inter-cell coordination, or hardware-specific power
measurement is attempted. These limitations are revisited in
Chapter~\ref{chap:oran-conclusion} in light of the simulation results
presented in Chapter~\ref{chap:oran-results}.

\section{Thesis Organization}
\label{sec:oran-organization}

The remainder of this thesis is organized as follows.
Chapter~\ref{chap:oran-literature} reviews the DRL and O-RAN
energy-optimization literature underlying this work, expanding on the gaps
identified in Section~\ref{sec:oran-problem-statement}.
Chapter~\ref{chap:oran-system-model} formulates the O-RAN energy-optimization
problem as a parameterized MDP and presents the BMPP-DQN architecture.
Chapter~\ref{chap:oran-results} presents the simulation setup, training
procedure, and results of the comparison against the DQN, DDPG, and MP-DQN
baselines. Chapter~\ref{chap:oran-conclusion} summarizes the thesis's
contributions, discusses the findings relative to the research questions
posed above, and outlines directions for future work.

% ---------------------------------------------------------------------------
% Local reference list for this chapter (IEEE numbered style). Entries below
% correspond to manuscript/ORAN_BMPP_DQN_Concept_Note_v1.md Section 9 and
% AGENTS.md's Foundational References table. Merge into a thesis-wide
% references.bib once main.tex exists; do not renumber these keys without
% also updating the \cite{} calls above.
% ---------------------------------------------------------------------------
\begin{thebibliography}{9}

\bibitem{qazzaz2026oreo}
M.~M.~H. Qazzaz, A. Salama, M. Hafeez, and S.~A.~R. Zaidi,
``OREO: Open RAN Energy Optimization via Deep Reinforcement Learning for 6G
Networks,'' \emph{IEEE Open Journal of the Communications Society}, vol.~7,
pp.~4165--4182, 2026.

\bibitem{oran2021mvp}
O-RAN Alliance, ``O-RAN Minimum Viable Plan and Commercialization
Strategy,'' Jun. 2021.

\bibitem{xiong2018pdqn}
J. Xiong, Q. Wang, Z. Yang, P. Sun, L. Han, Y. Zheng, H. Fu, T. Zhang, J. Liu,
and H. Liu, ``Parametrized Deep Q-Networks Learning: Reinforcement Learning
with Discrete-Continuous Hybrid Action Space,'' arXiv preprint
arXiv:1810.06394, 2018.

\bibitem{bester2019mpdqn}
C.~J. Bester, S.~D. James, and G.~D. Konidaris, ``Multi-Pass Q-Networks for
Deep Reinforcement Learning with Parameterised Action Spaces,'' arXiv
preprint arXiv:1905.04388, 2019.

\bibitem{tavakoli2018branching}
A. Tavakoli, F. Pardo, and P. Kormushev, ``Action Branching Architectures for
Deep Reinforcement Learning,'' in \emph{Proc. AAAI Conference on Artificial
Intelligence}, 2018.

\bibitem{bordin2025design}
M. Bordin, A. Lacava, M. Polese, S. Satish, M. AnanthaSwamy Nittoor,
R. Sivaraj, F. Cuomo, and T. Melodia, ``Design and Evaluation of Deep
Reinforcement Learning for Energy Saving in Open RAN,'' in \emph{Proc. IEEE
22nd Consumer Communications \& Networking Conference (CCNC)}, Las Vegas,
NV, USA, Jan. 2025, pp.~1--6.

\bibitem{sthankiya2024survey}
A. Sthankiya et al., ``A Survey on AI-Driven Energy Optimization in
Next-Generation RAN,'' arXiv preprint arXiv:2411.02164, 2024.

\bibitem{liang2026federated}
X. Liang, A. Al-Tahmeesschi, S.~B. Chetty, and H. Ahmadi, ``Green O-RAN
Operation: a Modern ML-Driven Network Energy Consumption Optimization,'' in
\emph{Proc. IEEE GLOBECOM}, 2025, pp.~4475--4480.

\bibitem{sohaib2024transfer}
M. Sohaib et al., ``DRL Transfer Learning for Cloud-Native O-RAN,''
arXiv preprint arXiv:2407.11563, 2024.

\end{thebibliography}
